% 근거: chapter4-prep-v1/A4/denominators.json,
%       chapter4-prep-v1/A6/observation-frequencies.jsonl,
%       chapter4-prep-v1/A6/proposed-three-candidates.json,
%       chapter4-prep-v1/A6/casebook.jsonl
\subsection{오류 유형 및 관측 사례}
\label{sec:ch4_errors}

Jev는 shortlist에서의 gold 탈락을 줄였으나, 이후 최종 선택되지 않는
gold도 남았다. 문항별 고유 gold 출현 443건의 최초 탈락 단계를 보면,
BM25 후보에는 있지만 shortlist에서 탈락한 gold는 수정 CE 42건,
Jev 13건이었다. shortlist에 남았으나 최종 선택되지 않은 gold는
각각 60건과 72건이었고, 최종 잔존은 233건과 250건이었다.
앞선 단계의 corpus 부재 33건과 공통 BM25 후보 부재 75건은
두 조건에서 같았다. 각 조건의 최초 탈락과 잔존을 합하면 443건이며,
이는 전역 고유 gold 386개와 다른 출현 단위 집계다.

이 단계별 집계는 exact gold의 회수ㆍ선택 상태를 설명한다.
shortlist 내 미선택에는 후보 목록당 하나를 고르는 제약과 중복 선택이
관련될 수 있어, 이를 곧바로 LLM의 법률 판단 오류로 볼 수는 없다.
gold 외 근거가 추가된 문항은 수정 CE 126개, Jev 116개였다.
gold 외 근거의 법률적 무관성은 이 집계로 판정할 수 없으며,
문항 수준의 관측 유형도 서로 겹치므로 빈도의 합은 전체 오류 문항
수가 아니다.

대표 사례는 Jev 개선, 수정 CE 우세, 지표 불일치 묶음에서
근거 F1 차이 절댓값의 중앙값에 가장 가까운 문항을 선정하고,
동률은 query\_id 사전순으로 정하였다.
표~\ref{tab:ch4_cases}는 이 규칙으로 고른 세 사례의 자동 점수를 보여준다.

\begin{table}[htbp]
    \centering
    \caption{대표 관측 사례의 자동 점수 변화 (수정 CE$\to$Jev)}
    \label{tab:ch4_cases}
    \small
    \setlength{\tabcolsep}{4pt}
    \begin{tabular}{@{}llrr@{}}
        \hline
        사례 & 질문 주제 & 근거 F1 & LF-Eval \\
        \hline
        A & 식품 보고 & $50.00\to100.00$ & $65.29\to99.18$ \\
        B & 복수 청구 & $80.00\to40.00$  & $71.79\to52.67$ \\
        C & 사면 효과 & $40.00\to0.00$   & $31.10\to30.55$ \\
        \hline
    \end{tabular}
    \par\smallskip
    \begin{minipage}{\linewidth}
        \footnotesize
        점수는 0--100 척도이다. 사례 식별자는 다음과 같다.\\
        A: \texttt{qa\_124\_2hop\_480\_rand}\\
        B: \texttt{qa\_324\_2hop\_515}\\
        C: \texttt{qa\_443\_3hop\_703}
    \end{minipage}
\end{table}

사례 A에서는 정답 근거의 완전 회수와 답변 문구의 변화가 함께
나타났다. 식품 생산실적과 위해식품 회수계획의 보고 상대방을 묻는
2-hop 문항으로, 수정 CE는 식품위생법 제42조제2항과 생산실적 관련
시행규칙을 선택하였다. Jev는 gold인 제42조제2항과 제45조제1항을
모두 선택하였으며, 답변에 회수계획의 보고 상대방으로
시장ㆍ군수ㆍ구청장을 포함하였다. 이 문항에서는 근거 F1과 LF-Eval이
모두 높아졌다.

사례 B에서는 Jev의 shortlist에서 gold가 탈락하고 두 자동 점수가
낮아졌다. 여러 당사자를 상대로 복수 청구를 제기하는 절차에 관한
2-hop 문항으로, 수정 CE는 gold인 민사소송법 제79조제1항과
제253조를 모두 선택하였다. Jev에서는 제79조제1항이 shortlist에서
탈락하고 최종 근거에 제83조제1항이 포함되었다. 답변은 각각
독립당사자참가와 공동소송참가 관련 내용을 담았으며, 근거 F1과
LF-Eval은 모두 수정 CE가 높았다. 각 제도의 실제 적용 가능성은
별도 법률 검토가 필요하다.

사례 C에서는 근거 F1과 답변 Token F1이 반대 방향으로 변했다.
사면과 복수 범죄의 처리에 관한 3-hop 문항으로, gold 3개 중 하나는
corpus에 정확히 매핑되지 않았고 다른 하나는 BM25 후보에 없었다.
생성된 후보 목록도 2개였다. 수정 CE는 gold 하나를 최종 선택했으나
Jev는 exact gold를 선택하지 않았다. 답변 Token F1은 18.60점에서
25.88점으로 높아져, 문자열 중복의 증가가 정답 근거 회수의 증가와
일치하지 않는 사례를 보여준다. 사면 유형, 범죄의 경합 관계,
gold 외 선택 조문의 타당성은 전문가 검토 사항이다.

세 사례는 법률 전문가의 판정이 완료되지 않은 자동 관측 사례이다.
모두 기존 사람 평가 60문항 표본 밖의 별도 정성 검토 대상으로,
블라인드 표본을 교체하거나 사람 평가 결과에 추가하지 않았다.
법률적 오류 여부와 제공 근거의 주장 지지는 미판정 상태이다.
